\documentclass[12pt,letterpaper]{article}
\usepackage[utf8]{inputenc}
\usepackage[spanish]{babel}
\usepackage[margin=2.5cm]{geometry}
\usepackage{amsmath}
\usepackage{amsfonts}
\usepackage{amssymb}

\begin{document}

\begin{center}
    \textbf{\Large Examen 1: Unidad I}\\
    \vspace{0.2cm}
    \textbf{Física para Ciencias de la Salud (005-1713)}\\
\end{center}

\vspace{0.5cm}
\noindent \textbf{Nombre y Apellido:} \underline{\hspace{7cm}} \hfill \textbf{C.I.:} \underline{\hspace{4cm}}

\vspace{0.5cm}
\noindent Resuelva los siguientes problemas, justificando claramente cada uno de sus procedimientos.

\vspace{0.5cm}

\begin{enumerate}

    \item \textbf{Ángulo plano:} Un paciente en rehabilitación realiza un movimiento de flexión con su brazo describiendo un ángulo de $75^\circ$. Exprese la medida de este ángulo en radianes.
    \vspace{3.5cm}

    \item \textbf{Sistemas de coordenadas:} En una placa de rayos X digitalizada 2D, el punto de origen de un hueso (articulación A) está en las coordenadas $(1, 2)$ cm y el extremo B en $(7, 10)$ cm. Calcule la longitud real del hueso.
    \vspace{3.5cm}
    
    \item \textbf{Vectores:} Sobre la vértebra L5 de una persona actúan dos fuerzas debido a la tensión muscular: $\vec{F}_1 = (45\hat{i} + 25\hat{j})$ N y $\vec{F}_2 = (-15\hat{i} + 35\hat{j})$ N. Encuentre el vector de la fuerza resultante que soporta la vértebra.
    \vspace{3.5cm}
    
    \item \textbf{Potencias de diez:} Un glóbulo blanco (leucocito) humano tiene un diámetro aproximado de $0.0000125$ metros. Exprese esta medida utilizando notación científica (potencias de diez) y el prefijo adecuado del Sistema Internacional.
    \vspace{3.5cm}

    \item \textbf{Sistemas de unidades:} En un análisis de laboratorio, se determina que la densidad del plasma sanguíneo es de $1.03 \, \text{g/cm}^3$. Convierta este valor a las unidades estándar del Sistema Internacional ($\text{kg/m}^3$).
    \vspace{3.5cm}

\end{enumerate}

\end{document}